% Generated draft: requires mathematical and source review.
\documentclass[11pt]{article}
\usepackage[margin=25mm]{geometry}
\usepackage{fontspec}
\setmainfont{Latin Modern Roman}
\usepackage{amsmath,amssymb,mathtools,unicode-math}
\setmathfont{Latin Modern Math}
\usepackage{enumitem,needspace,xurl,hyperref}
\setlength{\parindent}{0pt}
\setlength{\parskip}{5pt}
\setlength{\emergencystretch}{3em}
\newcommand{\N}{\mathbb{N}}
\newcommand{\Z}{\mathbb{Z}}
\newcommand{\Q}{\mathbb{Q}}
\newcommand{\R}{\mathbb{R}}
\newcommand{\C}{\mathbb{C}}
\newcommand{\F}{\mathbb{F}}
\newcommand{\A}{\mathbb{A}}
\newcommand{\PP}{\mathbb{P}}
\newcommand{\E}{\mathbb{E}}
\newcommand{\eps}{\varepsilon}
\newcommand{\dd}{\,\mathrm{d}}
\begin{document}
\section*{UnsolvedMath ID: 700842}
\textbf{Source problem code:} catalogue entry
\subsection{Definitions and mathematical statement}
For a group \(G\) and an integer \(k\geq 2\), let \(P_k(G)\) denote the property
\[
\forall (x_1,\ldots,x_k)\in G^k\;\exists\sigma\in\mathbb{S}_k\setminus\{\operatorname{id}\}:
\quad x_1\cdots x_k=x_{\sigma(1)}\cdots x_{\sigma(k)},
\]
where \(\mathbb{S}_k\) is the group of all permutations of \(\{1,\ldots,k\}\). Entries of the tuple may repeat, and the permutation may depend on the tuple. The index of permutability is
\[
p(G)=\min\{k\in\Z:k\geq 2\text{ and }P_k(G)\},
\]
provided this set is nonempty. The source does not specify a convention when the set is empty.

For a positive integer \(n\), let \(\mathbb{S}_n\) be the symmetric group on \(\{1,\ldots,n\}\). The problem proposed by M. Gutsan is to determine \(p(\mathbb{S}_n)\) as a function of \(n\). The source does not explicitly state the range of \(n\); positive finite degrees are the conventional interpretation used here, and inclusion of degree zero remains unspecified. The supplied registry labels the problem open; this label is not an independently verified mathematical or literature assessment.
\subsection{Short English statement}
Determine the least integer \(k\geq 2\) such that every ordered list of \(k\) elements of \(\mathbb{S}_n\), allowing repetitions, has the same product after some nonidentity permutation of its positions. The source leaves the range of \(n\) implicit.
\subsection{Sources}
\begin{itemize}
\item Supplied local UnsolvedMath registry record: permanent ID \(700842\), exact source problem code \texttt{catalogue entry}, title \emph{Kourovka Notebook Problem 12.23}; upstream revision \texttt{d371ee917ac9647fbb06cd6fc529deec6f3dda82}.
\item Supplied source attribution: \emph{The Kourovka Notebook (Unsolved Problems in Group Theory)}, 12th issue (1992), Problem 12.23; proposed by M. Gutsan.
\item The supplied background also identifies the 21st issue edition (2026), \texttt{arXiv:1401.0300v48}. The collection description names E. I. Khukhro and V. D. Mazurov as editors. These bibliographic details are supplied attribution and have not been independently checked.
\end{itemize}
\end{document}
